\setcounter{chaptercntr}{2}

\sectionbreak \section*{
  \gostTitleFont
  \redline
  \thechaptercntr .
  АНАЛИЗ СУЩЕСТВУЮЩИХ ТЕХНОЛОГИЧЕСКИХ РЕШЕНИЙ
}

\titlespace

\subsection*{
  \gostTitleFont
  \redline
  \thechaptercntr .\thesubchaptercntr \spc
  Обзор языков программирования и платформ для серверной разработки
} \addtocounter{subchaptercntr}{1}

\subtitlespace

{\gostFont

  \par \redline Разработка серверной части приложения такси относится к классу задач, для которых особенно важны производительность, устойчивость к отказам, расширяемость и поддержка большого числа одновременных запросов. Подобные требования характерны для высоконагруженных ИС, где ошибка в проектировании приводит не только к ухудшению качества обслуживания, но и к нарушению работы ключевых бизнес-процессов [2].

  \par \redline Для реализации серверной части могут использоваться различные языки программирования и платформы. Наиболее распространенными вариантами являются Java, Kotlin, C\#, JavaScript и Go. Каждый из них обладает собственными преимуществами, однако при проектировании корпоративных ИС важно учитывать зрелость экосистемы, наличие средств для работы с БД, поддержку тестирования, безопасность и возможности интеграции с внешними API.

  \par \redline Язык Java традиционно применяется для создания серверных приложений, поскольку он обладает устойчивой стандартной библиотекой, развитой экосистемой и большим количеством готовых решений для построения многослойной архитектуры. Дополнительным преимуществом является переносимость байт-кода между различными платформами и наличие зрелых средств разработки на уровне JDK и сопутствующих инструментов [4].

  \par \redline Среди серверных платформ Java особое место занимают решения, основанные на Spring. Данная экосистема позволяет строить приложения с четким разделением ответственности, использовать механизмы внедрения зависимостей, реализовывать REST API и упрощать работу с транзакциями, безопасностью и конфигурацией. Для учебного и прикладного проекта это снижает объем рутинного кода и делает структуру проекта более понятной [3].

  \par \redline В сравнении с Java, языки с более низким порогом входа, например JavaScript, удобны для быстрого прототипирования, но не всегда обеспечивают такой же уровень структурированности и зрелости в части enterprise-разработки. Go, в свою очередь, дает хорошие результаты при построении компактных высокопроизводительных сервисов, однако экосистема прикладных библиотек для типичных задач корпоративной ИС обычно менее разнообразна, чем у Java [2].

  \par \redline Отдельно следует отметить платформу .NET и язык C\#. Данное решение также широко используется для серверной разработки и обладает развитым набором средств для работы с web-приложениями, БД и механизмами безопасности. Однако в рамках данной работы использование Java выглядит более логичным, так как выбранная предметная область лучше соотносится с существующей экосистемой Spring и большим количеством типовых архитектурных решений, описанных в специализированной литературе.

  \par \redline Kotlin часто рассматривается как альтернатива Java благодаря лаконичному синтаксису и совместимости с JVM. Тем не менее при реализации серверной части приложения такси выбор Java является более оправданным, поскольку для него имеется больше типовых учебных примеров, стабильных практик проектирования и накопленного опыта использования в микросервисных системах.

  \par \redline При анализе платформ для серверной разработки также важно учитывать инструменты сборки, управления зависимостями и тестирования. Для Java-проектов такие задачи обычно решаются с помощью зрелых средств, которые хорошо интегрируются со Spring Boot и позволяют организовать воспроизводимую сборку, единый подход к конфигурации и автоматизацию проверок. Это особенно значимо для проекта, в котором несколько сервисов должны быть связаны общими соглашениями по API и едиными требованиями к качеству кода.

  \par \redline Таким образом, Java остается рациональным выбором для серверной части приложения такси. Она обеспечивает баланс между производительностью, сопровождаемостью и доступностью библиотек, а также хорошо сочетается с архитектурой, в которой отдельные функциональные блоки выделяются в независимые микросервисы [1].

}

\subtitlespace

\subsection*{
  \gostTitleFont
  \redline
  \thechaptercntr .\thesubchaptercntr \spc
  Выбор и обоснование стека технологий (Java, Spring Boot)
} \addtocounter{subchaptercntr}{1}

\subtitlespace

{\gostFont

  \par \redline При проектировании серверной части приложения такси необходимо не только выбрать язык программирования, но и определить набор технологий, который обеспечит быстрое создание прототипа, удобство сопровождения и возможность последующего расширения. В рамках данной работы в качестве основы выбран стек Java и Spring Boot, так как он позволяет реализовать типовую серверную ИС без избыточной сложности и с соблюдением распространенных архитектурных практик.

  \par \redline Spring Boot представляет собой подход к построению приложений на базе Spring, ориентированный на минимизацию конфигурации и ускорение запуска проекта. По документации платформы, Spring Boot упрощает создание автономных приложений, предоставляет готовые настройки для типичных сценариев и поддерживает построение сервисов, которые удобно разворачивать в контейнерной среде [5].

  \par \redline Для серверной части приложения такси это особенно важно, поскольку система должна включать обработку запросов пользователей, управление заказами, работу с авторизацией, интеграцию с внешними сервисами и обмен данными в формате JSON. Spring Boot предоставляет средства для организации контроллеров, сервисного слоя и доступа к данным, что соответствует слоистой структуре приложения и упрощает разделение ответственности между компонентами [5].

  \par \redline Еще одним преимуществом Spring Boot является удобство интеграции с типовыми инфраструктурными решениями. Для серверной части приложения это означает возможность подключать средства логирования, мониторинга, проверки состояния сервисов и внешние конфигурационные источники без существенного усложнения кода. В результате разработка концентрируется на предметной логике, а не на ручной сборке вспомогательных механизмов.

  \par \redline Существенным преимуществом Spring Boot является наличие встроенных механизмов для разработки REST API. Для приложения такси это означает возможность организовать обмен данными между клиентской и серверной частями в виде набора понятных HTTP-методов и ресурсов, что облегчает сопровождение интерфейсов и интеграцию с мобильными клиентами.

  \par \redline Отдельного внимания заслуживает поддержка безопасности. В ИС, где обрабатываются учетные данные пользователей и сведения о поездках, необходимо реализовать аутентификацию, авторизацию и защиту от несанкционированного доступа. Экосистема Spring предоставляет готовые средства, которые позволяют выстроить обработку JWT и контролировать доступ к защищенным ресурсам без ручной реализации базовых механизмов безопасности [5].

  \par \redline Практическая ценность Spring Boot проявляется и в том, что он хорошо поддерживает принцип разделения приложения на независимые слои. Контроллеры отвечают за прием и выдачу HTTP-запросов, сервисы содержат бизнес-логику, а уровень доступа к данным изолирует работу с БД. Такая структура соответствует принципам чистой архитектуры и упрощает последующее тестирование и сопровождение системы [4].

  \par \redline Еще одним аргументом в пользу выбранного стека является высокая распространенность Spring Boot в практике серверной разработки. Наличие большого числа документации, примеров и типовых решений снижает риск архитектурных ошибок и ускоряет выполнение проекта. Кроме того, подобный стек хорошо согласуется с микросервисным подходом, поскольку облегчает создание независимых сервисов с четко определенными API [1].

  \par \redline Таким образом, сочетание Java и Spring Boot обеспечивает достаточный уровень технологической зрелости для разработки серверной части приложения такси. Выбранный стек позволяет сосредоточиться на реализации предметной логики, не расходуя избыточные ресурсы на инфраструктурные детали, и при этом сохраняет возможности для дальнейшего развития системы [5].

}

\subtitlespace

\subsection*{
  \gostTitleFont
  \redline
  \thechaptercntr .\thesubchaptercntr \spc
  Анализ систем управления базами данных и средств кэширования
} \addtocounter{subchaptercntr}{1}

\subtitlespace

{\gostFont

  \par \redline Серверная часть приложения такси должна хранить данные о пользователях, заказах, маршрутах, статусах поездок, платежах и истории действий. Для таких задач требуется БД, которая обеспечивает целостность данных, поддержку транзакций и возможность выполнения сложных запросов. В большинстве случаев для этих целей применяются реляционные СУБД, поскольку именно они лучше всего подходят для работы с взаимосвязанными сущностями и строгими ограничениями целостности [6].

  \par \redline Среди реляционных СУБД особое место занимает PostgreSQL. Согласно документации, данная СУБД поддерживает транзакционность, индексацию, расширяемые типы данных и механизмы, необходимые для надежного хранения бизнес-данных [6]. Для приложения такси это важно, так как заказ, пользователь, водитель и платеж должны быть согласованы между собой и сохраняться без потери связности.

  \par \redline PostgreSQL также удобна тем, что позволяет применять сложные запросы и гибко организовывать схему данных. Это полезно при работе с отчетами, фильтрацией поездок, выборкой истории заказов и анализом активности пользователей. При проектировании серверной части подобная СУБД обеспечивает достаточную надежность и хорошо сочетается с ORM-подходом, который часто используется в Java-проектах.

  \par \redline Дополнительным преимуществом PostgreSQL является поддержка расширяемости и развитых механизмов индексации. Для сервиса такси это важно при обработке больших объемов данных, когда требуется быстро находить активные заказы, историю конкретного пользователя или ближайшие объекты по географическим признакам. Наличие таких возможностей делает PostgreSQL удобной основой для дальнейшего роста системы [6].

  \par \redline В качестве альтернативы для хранения данных можно рассматривать документные или key-value решения. Однако для приложения такси подобный выбор менее рационален, поскольку основная информация в системе тесно связана между собой и требует строгого контроля целостности. Именно поэтому реляционная модель данных оказывается предпочтительной для ядра серверной части.

  \par \redline Однако одной БД недостаточно для обеспечения высокой скорости отклика. В приложении такси часть данных используется значительно чаще, чем изменяется. К таким данным относятся справочники, настройки тарифов, сессионная информация, временные координаты и результаты часто повторяющихся проверок. Для снижения нагрузки на основную БД целесообразно применять кэширование.

  \par \redline Средства кэширования позволяют хранить часто запрашиваемые данные в оперативной памяти и тем самым уменьшать число обращений к основной СУБД. В качестве инструмента кэширования широко используется Redis. Документация Redis описывает поддержку различных структур данных и высокую скорость доступа, что делает данное решение практичным для сценариев, где важна малая задержка ответа [7].

  \par \redline Для сервиса такси Redis может быть полезен не только как обычный кэш, но и как средство хранения временных состояний, которые быстро меняются в ходе работы системы. Это относится к очередям на назначение заказа, к значениям, используемым при поиске ближайшего водителя, и к короткоживущим данным, которые нецелесообразно постоянно записывать в основную БД. Такое разделение снижает нагрузку на хранилище и повышает устойчивость сервиса к пиковым нагрузкам.

  \par \redline В приложении такси Redis может использоваться для хранения временных данных о состоянии заказа, ограничений на повторные запросы, промежуточных результатов вычислений и координат, обновляемых с высокой частотой. Это позволяет разгрузить PostgreSQL и повысить общую отзывчивость сервиса, особенно в периоды пикового спроса.

  \par \redline При сравнении БД и кэша следует учитывать, что они решают разные задачи. PostgreSQL обеспечивает долговременное и надежное хранение информации, тогда как Redis используется как быстрый слой доступа к данным. Следовательно, оптимальная схема для серверной части приложения такси предполагает совместное использование обеих технологий: основная БД отвечает за сохранность данных, а кэш - за скорость обработки наиболее частых запросов [6,7].

  \par \redline Таким образом, для разрабатываемой серверной части целесообразно использовать PostgreSQL в качестве основной БД и Redis в качестве средства кэширования. Такое сочетание обеспечивает баланс между надежностью хранения, производительностью и удобством последующего масштабирования серверной части приложения.

  \par
}

\setcounter{subchaptercntr}{1}
\setcounter{imagecntr}{1}
